%!TEX TS-program = xelatex

\documentclass[aspectratio=169,10pt]{beamer}

\newbool{russian}
\booltrue{russian}
\usepackage{HSE-theme/beamerthemeHSE}

% Русский язык и шрифты
\usepackage[english,russian]{babel}
\usepackage{fontspec}
\IfFontExistsTF{HSE Sans}{
  \setsansfont{HSE Sans}
  \setmonofont{Courier New}
}{
  % Резервный шрифт нужен только для систем, где HSE Sans не установлен.
  \setsansfont{DejaVu Sans}
  \setmonofont{DejaVu Sans Mono}
}
\usepackage{amsmath,amssymb,mathtools}
\uselanguage{russian}
\languagepath{russian}

\deftranslation[to=russian]{Theorem}{Теорема}
\deftranslation[to=russian]{Definition}{Определение}
\deftranslation[to=russian]{Example}{Пример}

\usepackage{booktabs}
\usepackage{tikz}
\usetikzlibrary{arrows.meta,positioning,calc}
\usepackage{pgfplots}
\pgfplotsset{compat=1.18}

% Небольшие вспомогательные настройки
\setbeamertemplate{blocks}[rounded][shadow=false]
\setbeamerfont{frametitle}{size=\Large}
\setbeamerfont{framesubtitle}{size=\small}
\setbeamerfont{block title}{size=\normalsize}
\setbeamerfont{block body}{size=\small}
\newcommand{\argmin}{\operatorname*{arg\,min}}
\newcommand{\argmax}{\operatorname*{arg\,max}}
\newcommand{\R}{\mathbb{R}}
\newcommand{\loss}{\ell}

% -----------------------------------------------------------------------------
% Система оценивания
% -----------------------------------------------------------------------------
\newcommand{\HWgrade}{H}
\newcommand{\CWgrade}{C}
\newcommand{\Colgrade}{K}
\newcommand{\Examgrade}{E}

\title[Методы оптимизации в ML]{Методы оптимизации\\в машинном обучении}
\subtitle{Лекция 1. Оптимизация как язык машинного обучения}
\author[Н. Лукьяненко]{Никита Лукьяненко}
\institute{Совместная программа СМОЛ ГУ и ФКН НИУ ВШЭ}
\date{2026/27 учебный год}

\begin{document}

% 1 ---------------------------------------------------------------------------
\frame[plain]{\titlepage}

% 2 ---------------------------------------------------------------------------
\begin{frame}{О преподавателе}
\begin{columns}[T,onlytextwidth]
  \column{0.58\textwidth}
  \begin{itemize}
    \item Никита Лукьяненко
    \item старший преподаватель ФКН НИУ ВШЭ
    \item преподавательский опыт: теория вероятностей, математическая статистика, статистическая теория обучения
    \item интересы курса: математическая сторона обучения моделей и вычислительное поведение алгоритмов
  \end{itemize}
  \vspace{0.5em}
  \begin{block}{Как будем работать}
    Моя цель --- не перечислить как можно больше методов, а научить вас понимать, \textbf{что именно оптимизируется, почему алгоритм устроен так и когда он может работать плохо}.
  \end{block}
  \column{0.36\textwidth}
  \begin{tikzpicture}[node distance=0.45cm]
    \node[draw=HSEblue,rounded corners,minimum width=4.5cm,minimum height=1cm,align=center] (m1) {Математика};
    \node[draw=HSEblue,rounded corners,below=of m1,minimum width=4.5cm,minimum height=1cm,align=center] (m2) {Алгоритмы};
    \node[draw=HSEblue,rounded corners,below=of m2,minimum width=4.5cm,minimum height=1cm,align=center] (m3) {Эксперименты};
    \draw[-{Latex},thick,HSEblue] (m1) -- (m2);
    \draw[-{Latex},thick,HSEblue] (m2) -- (m3);
  \end{tikzpicture}
\end{columns}
\end{frame}

% 3 ---------------------------------------------------------------------------
\begin{frame}{Что будет в курсе}
\framesubtitle{Пять вопросов, к которым мы будем возвращаться весь год}
\centering
\begin{tikzpicture}[node distance=0.52cm, every node/.style={align=center}]
  \node[draw=HSEblue,rounded corners,minimum width=8.5cm,minimum height=0.72cm] (q1) {Как формализовать задачу обучения?};
  \node[draw=HSEblue,rounded corners,below=of q1,minimum width=8.5cm,minimum height=0.72cm] (q2) {Как устроена геометрия целевой функции?};
  \node[draw=HSEblue,rounded corners,below=of q2,minimum width=8.5cm,minimum height=0.72cm] (q3) {Как найти минимум?};
  \node[draw=HSEblue,rounded corners,below=of q3,minimum width=8.5cm,minimum height=0.72cm] (q4) {Почему алгоритм сходится и насколько быстро?};
  \node[draw=HSEblue,rounded corners,below=of q4,minimum width=8.5cm,minimum height=0.72cm] (q5) {Что меняется в больших задачах машинного обучения?};
  \foreach \a/\b in {q1/q2,q2/q3,q3/q4,q4/q5}{\draw[-{Latex},HSEblue,thick] (\a)--(\b);}
\end{tikzpicture}
\vspace{0.7em}

\textbf{Главный герой первого семестра --- градиентный спуск.}
\end{frame}

% 4 ---------------------------------------------------------------------------
\begin{frame}{Как устроен курс}
\centering
\vspace{0.15em}

\begin{tikzpicture}[
  node distance=0.15cm,
  every node/.style={align=center},
  stage/.style={
    draw=HSEblue,
    rounded corners,
    text width=2.80cm,
    minimum width=3.00cm,
    minimum height=1.22cm,
    inner sep=3pt,
    font=\small
  }
]
  \node[stage] (lec) {\textbf{Лекция}\\идея и теория};
  \node[stage,right=of lec] (sem) {\textbf{Семинар}\\задачи руками};
  \node[stage,right=of sem] (lab) {\textbf{Лабораторная}\\эксперимент\\в Python};
  \node[stage,right=of lab] (ctrl) {\textbf{Контроль}\\самостоятельная\\работа};

  \draw[-{Latex},thick,HSEblue] (lec) -- (sem);
  \draw[-{Latex},thick,HSEblue] (sem) -- (lab);
  \draw[-{Latex},thick,HSEblue] (lab) -- (ctrl);
\end{tikzpicture}

\vspace{0.9em}
\begin{columns}[T,onlytextwidth]
\column{0.47\textwidth}
\begin{block}{На лекциях}
Определения, геометрическая интуиция, ключевые утверждения и небольшое число доказательств.
\end{block}

\column{0.47\textwidth}
\begin{block}{На практике}
Сначала маленькие задачи, где алгоритм можно увидеть целиком; затем --- модели и реальные данные.
\end{block}
\end{columns}

\vspace{0.35em}
\small
\textbf{В течение семестра:} 4 домашние лабораторные работы.
\end{frame}



% 5 ---------------------------------------------------------------------------
\begin{frame}{Система оценивания}
\small

\begin{center}
\begin{tabular}{@{}lcc@{}}
\toprule
\textbf{Активность} & \textbf{Обозначение} & \textbf{Вес} \\
\midrule
Домашние задания / 4 лабораторные работы & $H$ & $0.20$ \\
Контрольная работа & $C$ & $0.25$ \\
Коллоквиум & $K$ & $0.25$ \\
Экзамен & $E$ & $0.50$ \\
\bottomrule
\end{tabular}
\end{center}

\vspace{0.35em}
Каждая активность оценивается целым числом от $0$ до $10$.

\begin{block}{Расчёт}
% ВАЖНО: коэффициенты 0.20+0.25+0.25+0.50 суммируются в 1.20.
% Формула ниже сохранена ровно в том виде, который был задан для курса.
\[
S=0.20H+0.25C+0.25K+0.50E.
\]
Итоговая оценка получается арифметическим округлением рассчитанного результата.
\end{block}

\vspace{0.2em}
\textbf{Ориентировочные сроки:}
контрольная работа --- первая половина ноября; коллоквиум --- вторая половина ноября.

\vspace{0.45em}
\textbf{Домашние работы:}
4 лабораторные работы в течение семестра; точные дедлайны объявляются заранее.
\end{frame}



% 6 ---------------------------------------------------------------------------
\begin{frame}{Начнём с самой простой задачи}
\begin{columns}[T,onlytextwidth]
\column{0.47\textwidth}
\vspace{1em}
Рассмотрим функцию
\[
  f(x)=(x-3)^2+1.
\]

\begin{block}{Вопрос}
Какое значение $x$ нужно выбрать, чтобы сделать $f(x)$ как можно меньше?
\end{block}

\pause
\[
  x^\star=3, \qquad f(x^\star)=1.
\]

\column{0.49\textwidth}
\begin{tikzpicture}
\begin{axis}[
  width=7.3cm,height=5.4cm,
  axis lines=middle,
  xmin=-1,xmax=7,ymin=0,ymax=18,
  xlabel={$x$},ylabel={$f(x)$},
  samples=100,
  xtick={0,3,6},ytick={1,5,10,15},
  clip=false
]
\addplot[very thick,HSEblue,domain=-0.8:6.8] {(x-3)^2+1};
\addplot[only marks,mark=*,mark size=2.5pt,HSEred] coordinates {(3,1)};
\node[anchor=west] at (axis cs:3.2,1.3) {$x^\star$};
\end{axis}
\end{tikzpicture}
\end{columns}
\end{frame}

% 7 ---------------------------------------------------------------------------
\begin{frame}{Общая задача оптимизации}
\vspace{0.25em}

\begin{block}{Что означает запись?}
\centering
\[
\boxed{\displaystyle \min_{x\in X} f(x)}
\]
Нужно выбрать \textbf{допустимое} значение $x\in X$, при котором значение $f(x)$ как можно меньше.
\end{block}

\vspace{0.55em}
\begin{columns}[T,onlytextwidth]
\column{0.31\textwidth}
\centering
{\Large\color{HSEblue}$x$}\\[0.2em]
\textbf{что выбираем}\\[-0.1em]
\small переменная оптимизации

\column{0.31\textwidth}
\centering
{\Large\color{HSEblue}$f(x)$}\\[0.2em]
\textbf{что минимизируем}\\[-0.1em]
\small целевая функция

\column{0.31\textwidth}
\centering
{\Large\color{HSEblue}$X$}\\[0.2em]
\textbf{где можем выбирать}\\[-0.1em]
\small допустимое множество
\end{columns}

\vspace{0.9em}
\begin{alertblock}{Типичная интерпретация в машинном обучении}
\centering
$x$ --- параметры модели, \qquad $f(x)$ --- мера ошибки модели.
\end{alertblock}
\end{frame}


% 8 ---------------------------------------------------------------------------
\begin{frame}{Минимум и минимизатор --- разные объекты}
\begin{columns}[T,onlytextwidth]
\column{0.48\textwidth}
\begin{block}{Минимизатор}
Точка (или множество точек), где достигается наименьшее значение:
\[
  x^\star\in\argmin_x f(x).
\]
\end{block}
\column{0.48\textwidth}
\begin{block}{Минимальное значение}
Число --- значение функции в оптимальной точке:
\[
  f^\star=\min_x f(x).
\]
\end{block}
\end{columns}
\vspace{1em}
Для $f(x)=(x-3)^2+1$:
\[
  \argmin_x f(x)=\{3\},\qquad \min_x f(x)=1.
\]
\vfill
\centering
\Large\textbf{$\argmin$ = где? \qquad $\min$ = сколько?}
\end{frame}

% 9 ---------------------------------------------------------------------------
\begin{frame}{Минимизатор может быть не единственным}
\begin{columns}[T,onlytextwidth]
\column{0.47\textwidth}
\[
  f(x)=(x^2-1)^2.
\]
\[
  \argmin_x f(x)=\{-1,1\}.
\]

\begin{alertblock}{Важно}
$\argmin$ в общем случае --- \textbf{множество}, а не одна точка.
\end{alertblock}
\column{0.49\textwidth}
\begin{tikzpicture}
\begin{axis}[
  width=7.2cm,height=5.3cm,
  axis lines=middle,xmin=-2,xmax=2,ymin=-0.3,ymax=9,
  xlabel={$x$},ylabel={$f(x)$},samples=120,
  xtick={-1,0,1},ytick={0,4,8}
]
\addplot[very thick,HSEblue,domain=-1.9:1.9] {(x^2-1)^2};
\addplot[only marks,mark=*,mark size=2.5pt,HSEred] coordinates {(-1,0) (1,0)};
\end{axis}
\end{tikzpicture}
\end{columns}
\end{frame}

% 10 --------------------------------------------------------------------------
\begin{frame}{А минимум вообще всегда существует?}
\begin{columns}[T,onlytextwidth]
\column{0.31\textwidth}
\begin{block}{$f(x)=x^2$}
Минимум существует:
\[
\min f=0.
\]
\centering $x^\star=0$
\end{block}
\column{0.31\textwidth}
\begin{block}{$f(x)=e^x$}
\[
\inf f=0,
\]
но $e^x>0$ при любом конечном $x$.

\centering Минимума нет.
\end{block}
\column{0.31\textwidth}
\begin{block}{$f(x)=-x^2$}
\[
f(x)\to-\infty
\]
при $|x|\to\infty$.

\centering Функция не ограничена снизу.
\end{block}
\end{columns}
\vspace{1em}
\begin{alertblock}{Три фундаментальных вопроса курса}
\centering Существует ли решение? \quad Единственно ли оно? \quad Как его найти?
\end{alertblock}
\end{frame}

% 11 --------------------------------------------------------------------------
\begin{frame}{Добавление константы не меняет минимизаторы}
\begin{theorem}
Пусть $g(x)=f(x)+C$, где $C$ не зависит от $x$. Тогда
\[
  \argmin_x g(x)=\argmin_x f(x).
\]
\end{theorem}
\pause
\begin{proof}
Для любых $x_1,x_2$:
\[
 f(x_1)\le f(x_2)
 \quad\Longleftrightarrow\quad
 f(x_1)+C\le f(x_2)+C.
\]
Следовательно, добавление $C$ не меняет порядок значений функции и не меняет точки, в которых значение минимально.
\end{proof}
\end{frame}

% 12 --------------------------------------------------------------------------
\begin{frame}{Положительное масштабирование тоже не меняет минимизаторы}
\begin{theorem}
Если $g(x)=C f(x)$ и $C>0$, то
\[
  \argmin_x g(x)=\argmin_x f(x).
\]
\end{theorem}
\pause
Поскольку при $C>0$
\[
  f(x_1)\le f(x_2)
  \quad\Longleftrightarrow\quad
  Cf(x_1)\le Cf(x_2).
\]
\vspace{0.6em}
\begin{alertblock}{А если $C<0$?}
\[
  \argmin_x Cf(x)=\argmax_x f(x).
\]
Отрицательное масштабирование переворачивает порядок.
\end{alertblock}
\end{frame}

% 13 --------------------------------------------------------------------------
\begin{frame}{Почему это понадобится в машинном обучении}
Если $\loss_i(w)$ --- ошибка модели на объекте $i$, то часто встречаются две записи:
\[
  F_1(w)=\sum_{i=1}^{n}\loss_i(w),
  \qquad
  F_2(w)=\frac{1}{n}\sum_{i=1}^{n}\loss_i(w).
\]

Поскольку $1/n>0$,
\[
  \boxed{\argmin_w F_1(w)=\argmin_w F_2(w).}
\]

\begin{block}{То есть}
Для положения минимума сумма и среднее эквивалентны. Но их производные и численные масштабы отличаются --- позже это станет важным для выбора шага алгоритма.
\end{block}
\end{frame}

% 14 --------------------------------------------------------------------------
\begin{frame}{Не все преобразования безобидны}
\vspace{0.45em}

\begin{columns}[T,onlytextwidth]
\column{0.46\textwidth}
\begin{block}{Исходная задача}
Рассмотрим $x\in[-1,1]$ и
\[
f(x)=x.
\]
Тогда
\[
\argmin_{x\in[-1,1]} f(x)=\{-1\}.
\]
\end{block}

\column{0.46\textwidth}
\begin{block}{После преобразования}
Возведём целевую функцию в квадрат:
\[
g(x)=f(x)^2=x^2.
\]
Теперь
\[
\argmin_{x\in[-1,1]} g(x)=\{0\}.
\]
\end{block}
\end{columns}

\vspace{0.6em}
\begin{alertblock}{Вывод}
Даже простое преобразование целевой функции может изменить решение задачи оптимизации.
\end{alertblock}
\end{frame}


% 15 --------------------------------------------------------------------------
\begin{frame}{Небольшая задача}
Пусть
\[
  f(x)=(x-2)^2.
\]
Найдите минимизаторы следующих функций:
\[
\begin{aligned}
  &f(x),\\[0.35em]
  &3f(x),\\[0.35em]
  &f(x)+100,\\[0.35em]
  &-f(x).
\end{aligned}
\]
\vfill
\centering
\textbf{На обсуждение: 2 минуты.}
\end{frame}

% 16 --------------------------------------------------------------------------
\begin{frame}{Что делает модель машинного обучения?}
Пусть есть обучающая выборка
\[
  (x_i,y_i),\qquad i=1,\dots,n.
\]
Модель с параметрами $w$ строит предсказание
\[
  \widehat y_i=h_w(x_i).
\]

\begin{block}{Что мы можем менять?}
Данные $x_i,y_i$ уже зафиксированы. Во время обучения мы подбираем \textbf{параметры модели $w$}.
\end{block}

\[
  \boxed{\text{обучение} = \text{выбор хорошего значения }w}
\]
\end{frame}

% 17 --------------------------------------------------------------------------
\begin{frame}{Функция потерь: насколько плохо одно предсказание?}
\begin{definition}
Функция потерь $\loss(\widehat y,y)$ сопоставляет предсказанию и правильному ответу число, измеряющее ошибку.
\end{definition}

Для регрессии естественный пример:
\[
  \loss(\widehat y,y)=(\widehat y-y)^2.
\]

\begin{columns}[T,onlytextwidth]
\column{0.43\textwidth}
Если $\widehat y=y$, то
\[
  \loss=0.
\]
Чем дальше предсказание от ответа, тем больше штраф.
\column{0.53\textwidth}
\begin{tikzpicture}
\begin{axis}[width=7.2cm,height=4.0cm,axis lines=middle,xmin=-3,xmax=3,ymin=0,ymax=9.5,xlabel={$\widehat y-y$},ylabel={$\loss$},samples=80,xtick={-2,-1,0,1,2},ytick={0,4,8}]
\addplot[very thick,HSEblue,domain=-3:3] {x^2};
\end{axis}
\end{tikzpicture}
\end{columns}
\end{frame}

% 18 --------------------------------------------------------------------------
\begin{frame}{От одного объекта ко всей выборке}
Ошибка на объекте $i$:
\[
  \loss_i(w)=\loss\bigl(h_w(x_i),y_i\bigr).
\]
Средняя ошибка на обучающей выборке:
\[
  \boxed{
  F(w)=\frac{1}{n}\sum_{i=1}^{n}\loss_i(w)
  }
\]
называется \textbf{эмпирическим риском}.

\pause
\vspace{0.6em}
Обучение модели превращается в задачу
\[
  \boxed{
  w^\star\in\argmin_w F(w)
  }
\]
\end{frame}

% 19 --------------------------------------------------------------------------
\begin{frame}{Универсальная схема обучения}
\centering
\begin{tikzpicture}[node distance=0.38cm and 0.35cm,every node/.style={align=center}]
  \node[draw=HSEblue,rounded corners,minimum width=2.1cm,minimum height=0.9cm] (data) {данные\\$(x_i,y_i)$};
  \node[draw=HSEblue,rounded corners,right=of data,minimum width=2.1cm,minimum height=0.9cm] (model) {модель\\$h_w$};
  \node[draw=HSEblue,rounded corners,right=of model,minimum width=2.1cm,minimum height=0.9cm] (pred) {предсказание\\$\widehat y_i$};
  \node[draw=HSEblue,rounded corners,right=of pred,minimum width=2.1cm,minimum height=0.9cm] (loss) {ошибка\\$\loss_i(w)$};
  \node[draw=HSEblue,rounded corners,right=of loss,minimum width=2.1cm,minimum height=0.9cm] (risk) {риск\\$F(w)$};
  \foreach \a/\b in {data/model,model/pred,pred/loss,loss/risk}{\draw[-{Latex},HSEblue,thick] (\a)--(\b);}
\end{tikzpicture}

\vspace{1.4em}
\[
  \boxed{\displaystyle \min_w F(w)}
\]
\vspace{0.6em}
\begin{block}{Где начинается наш курс}
Архитектура модели и выбор функции потерь задают $F$. Дальше возникает вопрос: \textbf{как эффективно минимизировать эту функцию по параметрам $w$?}
\end{block}
\end{frame}

% 20 --------------------------------------------------------------------------
\begin{frame}{Разные модели --- одна оптимизационная структура}
\begin{columns}[T,onlytextwidth]
\column{0.31\textwidth}
\begin{block}{Линейная регрессия}
\[
\widehat y=w^\top x+b
\]
\[
\min_{w,b}F(w,b)
\]
\end{block}
\column{0.31\textwidth}
\begin{block}{Логистическая регрессия}
\[
p=\sigma(w^\top x+b)
\]
\[
\min_{w,b}F(w,b)
\]
\end{block}
\column{0.31\textwidth}
\begin{block}{Нейронная сеть}
\[
\widehat y=h_w(x)
\]
\[
\min_w F(w)
\]
\end{block}
\end{columns}
\vspace{1em}
\centering
Меняются модель и функция потерь, но остаётся вопрос
\[
\boxed{\text{как найти хорошие параметры?}}
\]
\end{frame}

% 21 --------------------------------------------------------------------------
\begin{frame}{Пример 1: линейная регрессия}
\begin{columns}[T,onlytextwidth]
\column{0.45\textwidth}
В одномерном случае модель имеет вид
\[
  \widehat y=wx+b.
\]
Нужно подобрать два параметра:
\[
  w\quad\text{и}\quad b.
\]
\begin{block}{Интерпретация}
$w$ отвечает за наклон прямой, $b$ --- за вертикальный сдвиг.
\end{block}
\column{0.51\textwidth}
\begin{tikzpicture}
\begin{axis}[width=7.6cm,height=5.1cm,axis lines=left,xmin=0,xmax=6,ymin=0,ymax=8,xlabel={$x$},ylabel={$y$},xtick=\empty,ytick=\empty]
\addplot[only marks,mark=*,mark size=2pt,HSEblue] coordinates {(0.7,1.4)(1.4,2.5)(2.0,2.6)(2.7,4.0)(3.4,4.2)(4.1,5.7)(4.8,5.5)(5.4,6.9)};
\addplot[very thick,HSEred,domain=0.4:5.7] {1+1.05*x};
\end{axis}
\end{tikzpicture}
\end{columns}
\end{frame}

% 22 --------------------------------------------------------------------------
\begin{frame}{Среднеквадратичная ошибка}
Для наблюдений $(x_i,y_i)$ и модели $\widehat y_i=wx_i+b$ определим
\[
  F(w,b)=\frac{1}{n}\sum_{i=1}^{n}(wx_i+b-y_i)^2.
\]
Тогда обучение линейной регрессии --- это
\[
\boxed{
  \min_{w,b}\frac{1}{n}\sum_{i=1}^{n}(wx_i+b-y_i)^2
}
\]

\begin{columns}[T,onlytextwidth]
\column{0.47\textwidth}
\begin{block}{Переменные оптимизации}
\centering $w,b$
\end{block}
\column{0.47\textwidth}
\begin{block}{Целевая функция}
\centering $F(w,b)$
\end{block}
\end{columns}
\vfill
\centering \textbf{Сегодня мы только ставим задачу. Градиент этой функции появится позже.}
\end{frame}

% 23 --------------------------------------------------------------------------
\begin{frame}{Два пространства, которые важно не путать}
\begin{columns}[T,onlytextwidth]
\column{0.46\textwidth}
\textbf{Пространство данных}

Каждая пара $(w,b)$ задаёт прямую
\[
  y=wx+b.
\]
Мы видим, насколько хорошо она проходит через наблюдения.
\column{0.50\textwidth}
\textbf{Пространство параметров}

Каждой паре $(w,b)$ соответствует число
\[
  F(w,b).
\]
Оптимизация происходит именно здесь:
\[
  (w,b)\longmapsto F(w,b).
\]
\end{columns}
\vspace{1.1em}
\begin{alertblock}{Ключевая мысль}
Модель живёт в пространстве данных, а алгоритм оптимизации двигается в пространстве параметров.
\end{alertblock}
\end{frame}

% 24 --------------------------------------------------------------------------
\begin{frame}{Пример 2: бинарная классификация}
Пусть
\[
  y\in\{0,1\}.
\]
Линейное выражение $w^\top x+b$ может принимать любые вещественные значения, поэтому само по себе не является вероятностью.

Используем логистическую функцию:
\[
  \boxed{
  \sigma(z)=\frac{1}{1+e^{-z}}
  }
\]
\[
  p(y=1\mid x)=\sigma(w^\top x+b).
\]
\vspace{0.4em}
\centering
$\sigma(z)\in(0,1)$ для любого $z\in\R$.
\end{frame}

% 25 --------------------------------------------------------------------------
\begin{frame}{Логистическая функция потерь}
Обозначим
\[
  p_i=\sigma(w^\top x_i+b).
\]
Одна из стандартных функций потерь для бинарной классификации:
\[
  F(w,b)=-\frac1n\sum_{i=1}^{n}
  \Bigl[y_i\log p_i+(1-y_i)\log(1-p_i)\Bigr].
\]

И снова обучение имеет вид
\[
  \boxed{\min_{w,b}F(w,b).}
\]

\begin{block}{Пока принимаем формулу как данность}
Её вывод из принципа максимального правдоподобия не нужен для сегодняшней цели. Нам важно увидеть оптимизационную структуру.
\end{block}
\end{frame}

% 26 --------------------------------------------------------------------------
\begin{frame}{Всегда ли нужно минимизировать только ошибку?}
Представьте две модели:
\begin{itemize}
  \item первая почти идеально подстраивается под обучающие данные, но очень сложна;
  \item вторая ошибается немного сильнее, зато устроена заметно проще.
\end{itemize}

\vspace{0.6em}
В машинном обучении часто нужен \textbf{компромисс} между качеством подгонки и сложностью модели.

\begin{alertblock}{Идея регуляризации}
Добавим к ошибке штраф за нежелательные значения параметров.
\end{alertblock}
\end{frame}

% 27 --------------------------------------------------------------------------
\begin{frame}{Регуляризованная задача}
\vspace{0.15em}

\centering
\[
\boxed{
\min_w
\left[
\underbrace{\frac1n\sum_{i=1}^{n}\loss_i(w)}_{\text{ошибка на данных}}
+
\lambda\,
\underbrace{R(w)}_{\text{штраф}}
\right]
}
\]

\vspace{0.35em}
\begin{columns}[T,onlytextwidth]
\column{0.30\textwidth}
\centering
{\Large\color{HSEblue}$w$}\\[0.15em]
\textbf{параметры модели}\\[-0.05em]
\small то, что мы подбираем

\column{0.30\textwidth}
\centering
{\Large\color{HSEblue}$R(w)$}\\[0.15em]
\textbf{регуляризатор}\\[-0.05em]
\small штраф за нежелательные параметры

\column{0.30\textwidth}
\centering
{\Large\color{HSEblue}$\lambda\ge 0$}\\[0.15em]
\textbf{сила штрафа}\\[-0.05em]
\small задаёт баланс двух частей
\end{columns}

\vspace{0.45em}
\raggedright
\textbf{\color{HSEblue}Пример:}\quad
$R(w)=\|w\|_2^2$.

\vspace{0.25em}
\small Чем больше $\lambda$, тем сильнее целевая функция штрафует большие значения параметров.
\end{frame}



% 28 --------------------------------------------------------------------------
\begin{frame}{Оптимизация как компромисс}
\centering
\vspace{0.6em}
\Large
\[
\underbrace{\text{ошибка на данных}}_{\text{хочется уменьшить}}
\quad+
\lambda
\underbrace{\text{сложность модели}}_{\text{хочется контролировать}}
\]
\normalsize
\vspace{1em}
\begin{block}{Важно}
Оптимизация не обязательно означает «минимизировать ошибку любой ценой». Мы сначала \textbf{конструируем целевую функцию}, которая отражает нужный нам компромисс, а уже потом ищем её минимум.
\end{block}
\end{frame}

% 29 --------------------------------------------------------------------------
\begin{frame}{Вся первая лекция --- в одной формуле}
\centering
\vspace{0.5em}
\[
\boxed{
  w^\star\in\argmin_w
  \left[
  \frac1n\sum_{i=1}^{n}
  \loss\bigl(h_w(x_i),y_i\bigr)
  +\lambda R(w)
  \right]
}
\]
\vspace{1em}
\begin{itemize}
  \item $w$ --- параметры модели;
  \item $h_w$ --- модель;
  \item $\loss$ --- ошибка на одном объекте;
  \item $R$ --- регуляризатор;
  \item $\lambda$ --- сила регуляризации.
\end{itemize}
\end{frame}

% 30 --------------------------------------------------------------------------
\begin{frame}{Но как найти $w^\star$?}
\centering
\Large
Если параметров два, поверхность функции можно хотя бы нарисовать.
\vspace{0.7em}

А если параметров --- миллион?
\vspace{1em}

\begin{alertblock}{Следующий вопрос курса}
\centering
Как использовать геометрию функции, чтобы понять, \textbf{куда двигаться из текущей точки}?
\end{alertblock}
\vspace{0.7em}
\normalsize
Следующая лекция: \textbf{линейная алгебра для оптимизации}.
\end{frame}

% 31 --------------------------------------------------------------------------
\begin{frame}{Что нужно унести с собой}
\begin{enumerate}
  \item Задача оптимизации имеет вид $\min_{x\in X}f(x)$.
  \item $\argmin$ обозначает точки-минимизаторы, $\min$ --- минимальное значение.
  \item Обучение модели часто сводится к минимизации среднего значения функции потерь.
  \item Регуляризация добавляет в целевую функцию штраф и меняет компромисс, который мы оптимизируем.
  \item Следующая задача --- научиться находить минимум алгоритмически.
\end{enumerate}
\vfill
\centering
\textbf{К следующему занятию: повторить векторы, матрицы, скалярное произведение и собственные значения.}
\end{frame}

\end{document}